\documentclass[11pt]{article}

% Template-agnostic preamble --- works without any venue style file.
\usepackage[margin=1in]{geometry}
\usepackage[utf8]{inputenc}
\usepackage[T1]{fontenc}
\usepackage{hyperref}
\hypersetup{colorlinks=true, linkcolor=blue, citecolor=blue, urlcolor=blue,
  pdftitle={Multi-Channel Mechanistic Signatures of Agent WANDERING},
  pdfauthor={Caio Vicentino}}
% Overflow protection for long HF/dataset paths (per LaTeX long-path rule)
\emergencystretch=3em
\sloppy
\hyphenpenalty=2000
\exhyphenpenalty=0
\hbadness=10000
\usepackage{booktabs}
\usepackage{amsfonts}
\usepackage{amsmath}
\usepackage{amssymb}
\usepackage{xcolor}
\usepackage{graphicx}
\usepackage{caption}
\usepackage{tabularx}
\usepackage{array}
\usepackage{multirow}

\title{Multi-Channel Mechanistic Signatures of Agent WANDERING:\\
Classification, Causal Localization, and Behavior-Legible Response to Intervention}

\author{%
  Caio Vicentino \\
  OpenInterpretability \\
  Fortaleza, Brazil \\
  ORCID: 0009-0003-4331-6259 \\
  \texttt{caio@openinterp.org}%
}
\date{}

\begin{document}
\maketitle

\begin{abstract}
WANDERING --- the failure mode in which an LLM agent stays internally confident it has solved a task yet never emits a termination action and exhausts its turn budget --- is a 34\% blind spot in probe-based agent monitoring \cite{toolentropy2026}. We characterize WANDERING mechanistically on $N{=}99$ Qwen3.6-27B SWE-bench Pro trajectories (SUCCESS=40, LOCKED=39, WANDERING=20) by extracting 60 multi-channel features (text, tool-use, per-layer residual at L11/23/31/43/55, and temporal) and fitting a 3-way L1 classifier under nested cross-validation. After removing nine label-leaking features we report a clean macro-F1 of $0.636\pm0.035$ ($z{=}{+}5.88$, $p{=}0.001$ vs.\ a 1000-permutation null). SUCCESS is robustly separable (AUROC $1.0$, driven by tool diversity), but the LOCKED--WANDERING boundary is mechanistically weak ($p{=}0.035$), echoing a text-based human taxonomy \cite{liu2026failures} that places $\approx$60\% of both into a single ``Non-Progressive Iteration'' category. Stability selection independently surfaces the mid-to-edge mechanism hypothesized in \cite{toolentropy2026}: LOCKED is dominated by L43 mid-layer convergence, WANDERING by L11 edge-layer drift, with effect sizes that are statistically robust but small ($2$--$3$ percentage points in cosine). Finally, we ask whether the signature predicts \emph{which WANDERING agents flip} to \texttt{finish\_tool} during an L11 SUCCESS-direction injection run (paper~II). Paper~II finds the injection's \emph{direction} is not a causal rescue lever (paired McNemar $p{=}0.73$ against a same-session no-hook baseline that itself flips $7/20$ under run-instability); we therefore treat the $6/20$ observed flips as a behavioral-response label, not a causal-rescue label. The residual signature does not predict which agents flip (LOO-AUC $0.619$, $p{=}0.16$), but tool-entropy collapse depth does (AUC $0.768$, $p\approx0.06$): the more an agent has collapsed onto a single repeated tool, the more likely it tips over to a finish action. Response to intervention is residual-blind but behavior-legible. We release all code, features, and labels.
\end{abstract}

\section{Introduction}

Mechanistic interpretability for \emph{agents} --- as opposed to single forward passes --- is nascent. Long-horizon LLM agents fail in ways that are invisible to a single-token probe: the model may consolidate an internal ``I am done'' verdict mid-trajectory yet never translate it into a \texttt{finish\_tool} action, looping on its action space until it exhausts the turn budget. We call this \textbf{WANDERING}, following the operational definition of \cite{toolentropy2026}: a failed trajectory ($\texttt{finish\_reason}{=}\texttt{max\_turns}$) whose capability probe stays high (median final score $1.0$) --- the agent ``knows'' it has solved the task but never acts on that knowledge.

Companion work \cite{toolentropy2026} showed WANDERING is \emph{detectable} --- tool-entropy collapse over the last ten turns separates WANDERING from SUCCESS across three model families (Qwen3.6-27B, Llama-70B, GPT-5) --- and proposed a mechanism (\S13 therein): the mid-layer verdict is consolidated but an edge-layer (L11/L55) alignment step fails to emit the action. That paper, however, treated WANDERING as a single phenomenon detected by a single channel. Three questions remain open: (i) Is WANDERING mechanistically \emph{one} thing, or does it decompose? (ii) Does the proposed mid-to-edge mechanism survive an analysis that does not assume it? (iii) Is the signature merely descriptive, or does it predict which agents \emph{respond} behaviorally to an intervention?

We answer all three on the same $N{=}99$ corpus, with four contributions:

\begin{enumerate}
\item \textbf{A multi-channel signature.} 60 features across text, tool-use, residual, and temporal channels classify the three sub-classes (SUCCESS/LOCKED/WANDERING) at macro-F1 $0.636$ ($p{=}0.001$). We report a transparent walk-back from a leaky $0.987$ baseline.
\item \textbf{Independent mechanism confirmation.} Stability selection --- using none of \cite{toolentropy2026}'s detector designs --- recovers its mid-to-edge picture: LOCKED $\to$ L43 mid-layer freeze, WANDERING $\to$ L11 edge-layer drift. Effect sizes are robust but small, and we say so.
\item \textbf{A bridge to a human taxonomy.} An LLM-judge mapping to \cite{liu2026failures} places $\approx$60\% of both LOCKED and WANDERING into one text category, matching our finding that the LOCKED--WANDERING boundary is mechanistically weak ($p{=}0.035$).
\item \textbf{Behavior-legible response to intervention.} Paper~II finds the L11 SUCCESS-direction injection is not a causal rescue lever (paired McNemar $p{=}0.73$); using the $6/20$ observed \texttt{finish\_tool} flips as a behavioral-response label, we show the residual signature does \emph{not} predict which agents flip (LOO-AUC $0.619$, $p{=}0.16$) but tool-entropy collapse depth does (AUC $0.768$, $p\approx0.06$) --- the intervention-response signal is behavioral, not residual.
\end{enumerate}

\section{Related Work}

\textbf{Agent failure taxonomies.} \cite{liu2026failures} build a 3-phase / 9-category / 20-subcategory taxonomy of SWE-bench failures via human open-coding (Cohen's $\kappa{=}0.72$--$0.77$); their ``Non-Progressive Iteration'' (C2.1) and ``Context Amnesia'' (C3) overlap conceptually with WANDERING but are not automatically detectable from model internals. \cite{majgaonkar2026codeagent} characterize success/failure code-agent trajectories with length-distribution statistics (``Quick'' vs.\ ``Extended'' failures), text-only and without a formal taxonomy. \cite{whenagents2026} course-correct SWE agents with inference-time process reward models, a \emph{text}-level intervention with no model-internal signal. None use residual-stream features or predict which agents respond to an intervention.

\textbf{Mechanistic interpretability of agent tool use.} \cite{beyondblackbox2026} apply sparse autoencoders and linear probes to pre-action residuals to predict tool-need and tool-risk, but evaluate on function-calling benchmarks rather than long-horizon failure, and do not cluster or causally test failure modes. Our work is, to our knowledge, the first to combine residual-stream features, multi-channel classification, a human-taxonomy bridge, and intervention-response prediction on the same agent-trajectory corpus.

\textbf{Detection of WANDERING.} The companion paper \cite{toolentropy2026} establishes the detectors (forensic text $v1$, cross-layer probe disagreement $v4$ with a $\sim$15-turn early-warning lead, tool-entropy $v5$) and their cross-architecture generalization, plus a cross-task null on METR MALT ($p{=}0.81$) scoping the claim to multi-turn code-execution agents. The present paper takes those sub-class labels as given and asks what the internals say.

\section{Data and Methods}

\textbf{Corpus.} $N{=}99$ multi-turn trajectories from Qwen3.6-27B (\texttt{enable\_thinking=True}, temperature $1.0$) on SWE-bench Pro (qutebrowser / openlibrary / ansible), collected on NVIDIA RTX 6000 Pro Blackwell. Per-turn residual activations are captured at L11/23/31/43/55 (\texttt{bfloat16}, $d_{\text{model}}{=}5120$). The dataset (99 traces + 99 patches + per-turn residual safetensors + derived features) is public:
\begin{center}\small\texttt{huggingface.co/datasets/caiovicentino1/swebench-pro-qwen36-27b-phase6}\end{center}

\textbf{Sub-class derivation.} SUCCESS: test passes after \texttt{finish\_tool} ($n{=}40$). LOCKED: failed and probe stably $<0.40$ ($n{=}39$). WANDERING: failed and probe stays $>0.70$ with no \texttt{finish\_tool} ($n{=}20$).

\textbf{Feature channels (63 raw $\to$ 60 used).} \emph{Text} (11): doubt / exploration / completion word densities, thinking and content lengths, n-gram overlap. \emph{Tool-use} (9): Shannon entropy of tool calls (full / first-10 / last-10), unique-tool count, top-tool fraction, switching rate, longest same-tool run, calls per turn. \emph{Residual} (32): per-layer norm mean/std/late-mean, consecutive-turn cosine (mean and late), first-vs-last drift cosine, plus two cross-layer aggregates. \emph{Temporal} (11): turn count, patch size, probe-trajectory statistics, detector fire-turns.

\textbf{Label-leakage control.} Nine features are removed before the clean analysis: \texttt{probe\_score\_*} (six features derived from the same probe used to bin sub-class via the lock threshold), \texttt{n\_turns} and \texttt{patch\_n\_bytes} (termination proxies), and \texttt{detector\_v4\_fire\_turn} (assignment-process correlate). Two all-NaN detector columns and one zero-variance column are also dropped, leaving 52 features for the primary model.

\textbf{Pipeline.} 3-way L1 multinomial logistic regression (saga), nested 5$\times$3 cross-validation (inner GridSearch over $C\in\{10^{-3},\dots,10^2\}$; modal $C{=}1.0$, interior to the grid), \texttt{StandardScaler} inside the pipeline (no test leakage), 1000-permutation null with the Phipson--Smyth $(1{+}\#\{\text{null}\ge\text{obs}\})/(1{+}n)$ $p$-value, and stability selection over 200 bootstraps (feature stable iff selected in $\ge 80\%$, Meinshausen--B\"uhlmann).

\section{Results}

\subsection{Three-way classification, with a transparent walk-back}

We report the leaky baseline alongside the cleaned analysis (per a pre-registered transparency standard):

\begin{center}
\begin{tabular}{lccc}
\toprule
Run & Macro-F1 & $z$ & $p$ \\
\midrule
Leaky baseline (uses probe / termination proxies) & $0.987\pm0.030$ & $+12.43$ & $0.001$ \\
\textbf{Clean (primary)} & $\mathbf{0.636\pm0.035}$ & $\mathbf{+5.88}$ & $\mathbf{0.001}$ \\
\bottomrule
\end{tabular}
\end{center}

The leaky model is near-perfect only because \texttt{probe\_score\_*} is definitionally tied to the sub-class boundary; removing it drops F1 to $0.636$, still highly significant against the permutation null (mean $0.322\pm0.053$). Per-class AUROC: SUCCESS $1.000$, LOCKED $0.864$, WANDERING $0.795$. The confusion matrix shows the structure plainly:
\begin{center}
\begin{tabular}{lccc}
\toprule
true $\backslash$ pred & LOCKED & SUCCESS & WANDERING \\
\midrule
LOCKED & 24 & 0 & 15 \\
SUCCESS & 0 & 40 & 0 \\
WANDERING & 13 & 0 & 7 \\
\bottomrule
\end{tabular}
\end{center}
SUCCESS is perfectly separated; the 28 LOCKED$\leftrightarrow$WANDERING confusions are the entire error budget.

\textbf{Pairwise contrasts} confirm: WANDERING-vs-SUCCESS and SUCCESS-vs-LOCKED both reach macro-F1 $1.000$ ($p{=}0.005$), but \textbf{WANDERING-vs-LOCKED is only $0.631$ ($z{=}1.81$, $p{=}0.035$)} --- barely significant. The two failure modes share mechanistic territory.

\subsection{Stability selection recovers the mid-to-edge mechanism}

Three features are stable ($\ge 0.80$ selection frequency), one per class:

\begin{center}
\begin{tabular}{llcc}
\toprule
Feature & Class & Sel.\ freq & Per-class median (S / L / W) \\
\midrule
\texttt{tool\_diversity\_count} & SUCCESS & $1.00$ & $3.0$ / $2.0$ / $2.0$ \\
\texttt{L43\_cosine\_consec\_late} & LOCKED & $0.835$ & $0.769$ / $\mathbf{0.839}$ / $0.814$ \\
\texttt{L11\_drift\_first\_last} & WANDERING & $0.95$ & $0.973$ / $0.975$ / $\mathbf{0.956}$ \\
\bottomrule
\end{tabular}
\end{center}

This independently reproduces the \cite{toolentropy2026} \S13 picture, using none of its detector designs: \textbf{LOCKED} is marked by high L43 mid-layer consecutive-turn cosine (the mid-layer representation \emph{freezes} late --- a consolidated wrong verdict that stops updating); \textbf{WANDERING} by low L11 first-vs-last drift cosine (the edge-layer representation keeps \emph{rotating} --- the model never consolidates to an action); \textbf{SUCCESS} by tool diversity (it reaches the third available tool). An ultra-conservative re-fit that drops \emph{all} L43 features and \texttt{tool\_diversity} still surfaces \texttt{L11\_drift\_first\_last} in the top three for the WANDERING-vs-LOCKED contrast (AUC $0.696$, $z{=}1.88$, $p{=}0.030$), so the edge-layer signal does not depend on the mid-layer or tool features.

\textbf{Effect sizes are subtle.} The discriminating medians differ by only $2$--$3$ percentage points on a $0$--$1$ cosine scale, with overlapping inter-quartile ranges. The signal is a statistical regularity captured by L1 selection across folds, not a visible gap; we make no stronger claim.

\subsection{Bridge to a human taxonomy}

Mapping all 99 trajectories to \cite{liu2026failures} via an LLM judge: SUCCESS maps cleanly to a ``no-failure'' sentinel (39/40), but \textbf{60\% of both LOCKED (23/39) and WANDERING (12/20) map to the single category C2.1 ``Non-Progressive Iteration.''} A text-based taxonomy cannot separate the two sub-types, exactly matching the mechanistic weakness of the LOCKED--WANDERING boundary ($p{=}0.035$). The probe-derived sub-class definitions may be imposing categorical boundaries on what is, behaviorally and mechanistically, a continuum.

\subsection{Behavior-legible response to intervention}

Paper~II finds L11 SUCCESS-direction injection is \emph{not} a causal rescue lever: paired against a same-session no-hook baseline that itself flips $7/20$ under run-instability, the injection's effect is null (McNemar $p{=}0.73$). Nonetheless $6/20$ WANDERING agents emitted \texttt{finish\_tool} during the injection run. We ask, \emph{pre-registered}, a flip-proneness question --- which agents flip, independent of whether the injection's \emph{direction} causes it: does the multi-channel signature predict which agents flip, using only features from the original no-hook trajectory? Decision gate: GO iff leave-one-out AUC $\ge 0.70$ and permutation $p < 0.10$.

\textbf{The pre-registered mechanistic model fails the gate} (LOO-AUC $0.619$, $p{=}0.16$). The residual channels that classify the sub-classes do not predict which agents flip. But a single \emph{behavioral} channel does (exploratory, flagged as not the pre-registered primary):

\begin{center}
\begin{tabular}{lccc}
\toprule
Feature & Univariate AUC & Perm $p$ & $\Delta$ median (resc.\ $-$ resist.) \\
\midrule
\texttt{L11\_drift\_first\_last} & $0.560$ & --- & $-0.010$ \\
\texttt{L11\_cosine\_consec\_late} & $0.619$ & --- & $+0.003$ \\
\texttt{L43\_cosine\_consec\_late} & $0.500$ & --- & $+0.001$ \\
\textbf{\texttt{tool\_entropy\_last10}} & $\mathbf{0.768}$ & $\mathbf{0.066}$ & $\mathbf{-0.487}$ \\
\bottomrule
\end{tabular}
\end{center}

Flipping agents had much \emph{lower} original tool-entropy (median $0.234$ vs.\ $0.722$; Mann--Whitney $p{=}0.060$): \textbf{the more an agent has collapsed onto a single repeated tool, the more likely it flips to \texttt{finish\_tool} under the injection run.} Agents in a tight tool-loop are near the finish threshold and tip over; agents flailing across many tools do not. (Whether the \emph{direction} of the injection causes the flip is null --- paper~II; tool-entropy predicts flip-proneness, not a causal rescue.) Strikingly, the \emph{same} tool-entropy signal that \emph{detects} WANDERING \cite{toolentropy2026} also predicts which agents respond behaviorally to an intervention --- the actionable information is behavioral, not residual. We treat this as exploratory ($N{=}20$, 6 events) and flag confirmatory replication as required; it nonetheless points to a tool-entropy-gated detection$\to$intervention loop as the natural next experiment.

\subsection{Run-to-run stability prediction}

Re-running the 20 WANDERING instances no-hook on the same hardware yields variable outcomes under temperature $1.0$ (an independent re-run flips 7/20 to \texttt{finish\_tool}; a fresh determinism check flips 0/5). Predicting which instances flip from their original features gives LOO-AUC $0.725$ ($p{=}0.096$, marginal), with \texttt{L11\_cosine\_consec\_late} the top predictor --- a second, independent appearance of an L11 edge-layer signal. This is a run-to-run (temperature-sampling) effect, \emph{not} a hardware effect: all runs used the same GPU.

\section{Discussion}

\textbf{WANDERING is a weak boundary, not a clean type.} Three independent axes agree: the supervised classifier ($p{=}0.035$), the human-taxonomy bridge ($\approx$60\% co-located), and the intervention-response analysis (heterogeneous) all say LOCKED and WANDERING share territory. The probe-threshold definition draws a line through a continuum. This is a caution for any agent-failure benchmark that treats such labels as ground truth.

\textbf{The mechanism survives an agnostic analysis --- subtly.} That stability selection, with no knowledge of \cite{toolentropy2026}'s mid-to-edge hypothesis, recovers L43-mid-for-LOCKED and L11-edge-for-WANDERING is genuine independent corroboration. But the $2$--$3$pp effect sizes mean this is a statistical signature, not a smoking gun; we resist over-claiming a clean mechanistic dissociation.

\textbf{Intervention-response is behavioral.} The most consequential result is the dissociation between \emph{descriptive} and \emph{intervention-response} signal: the residual signature distinguishes the sub-classes but does not predict which agents flip under an intervention, whereas tool-entropy --- a cheap behavioral statistic --- does. For a deployed monitor, this is good news: the cheapest channel is the one that indexes who responds to intervention, and it streams online. (The direction of the residual intervention itself is a null --- paper~II --- so this is a behavior-legible response signal, not a causal-rescue selector.)

\section{Limitations}

(i) Single model (Qwen3.6-27B) and single benchmark (SWE-bench Pro); the multi-channel phenotype is not yet cross-validated off-Qwen. (ii) $N{=}20$ WANDERING binds every within-WANDERING analysis; the intervention-response prediction ($6$ events) and run-stability prediction are underpowered and the tool-entropy result is exploratory. (iii) Effect sizes are small ($2$--$3$pp cosine). (iv) \texttt{L11\_drift\_first\_last} correlates with thinking-length ($\rho{=}0.63$); a partial-correlation control finds L11 retains signal ($\rho_{\text{partial}}{=}-0.29$, $p{=}0.004$) but a matched-subset test is marginal ($p{=}0.08$), so we cannot fully rule out a thinking-volume confound. (v) The detector implementations used for feature extraction differ slightly from \cite{toolentropy2026}'s; two detector-fire columns were all-NaN and dropped.

\section{Conclusion}

WANDERING is detectable and mechanistically a weak boundary rather than a clean type. Its residual signature confirms a mid-to-edge picture but does not tell you which agents respond to an intervention; the behavioral tool-entropy channel does (while the residual intervention's direction is itself a null, per paper~II). The detection signal and the response-selection signal are the same cheap statistic --- which is what makes a closed detection$\to$intervention loop worth testing, and is our next pre-registered target.

\paragraph{Reproducibility.} All code, the 60-feature matrix, sub-class and intervention-response labels, permutation nulls, and the LLM-judge prompt are released under Apache-2.0 at \texttt{github.com/OpenInterpretability/openinterp-swebench-harness}; the dataset (traces + residuals) is on HuggingFace at the link above. Companion papers: Tool-Entropy Collapse (DOI \texttt{10.5281/zenodo.20368601}); The Right Locus Is Still Not a Rescue Lever (in preparation).

\begin{thebibliography}{9}
\bibitem{toolentropy2026} C.\ Vicentino. \emph{Tool-Entropy Collapse: A Cross-Architecture Signature of Agent WANDERING Failure}. Zenodo, 2026. DOI 10.5281/zenodo.20368601.
\bibitem{liu2026failures} Liu et al. \emph{An Empirical Study on Failures in Automated Issue Solving}. arXiv:2509.13941, 2025.
\bibitem{majgaonkar2026codeagent} Majgaonkar et al. \emph{Understanding Code Agent Behaviour: An Empirical Study of Success and Failure Trajectories}. ICSE 2026 (arXiv:2511.00197).
\bibitem{whenagents2026} S.\ Gandhi et al. \emph{When Agents go Astray: Course-Correcting SWE Agents with PRMs}. NeurIPS 2025 (arXiv:2509.02360).
\bibitem{beyondblackbox2026} \emph{Beyond the Black Box: Interpretability of Agentic AI Tool Use}. arXiv:2605.06890, 2026.
\end{thebibliography}

\end{document}
